\documentclass[12pt,a4paper]{article}

\usepackage{fontspec}
\usepackage{polyglossia}
\setmainlanguage{french}

\usepackage[top=2.5cm, bottom=2.5cm, left=3cm, right=2.5cm]{geometry}
\usepackage{setspace}
\onehalfspacing

\usepackage{titlesec}
\titleformat{\section}{\large\bfseries}{}{0em}{\thesection\quad}
\titleformat{\subsection}{\normalsize\bfseries}{}{0em}{\thesubsection\quad}
\titleformat{\subsubsection}{\normalsize\itshape}{}{0em}{\thesubsubsection\quad}

\usepackage{booktabs}
\usepackage{array}
\usepackage{tabularx}
\usepackage{multirow}
\usepackage{xcolor}
\usepackage{mdframed}
\usepackage{listings}
\usepackage{enumitem}
\usepackage[hidelinks]{hyperref}
\usepackage{fancyhdr}

\definecolor{bleuCMR}{RGB}{0, 62, 126}
\definecolor{vertOK}{RGB}{0, 130, 60}
\definecolor{rougeERR}{RGB}{180, 30, 30}
\definecolor{orangeWARN}{RGB}{200, 100, 0}
\definecolor{grisLight}{RGB}{245, 245, 245}
\definecolor{grisCode}{RGB}{250, 250, 250}

\mdfdefinestyle{encadre}{
  backgroundcolor=grisLight,
  linecolor=bleuCMR,
  linewidth=1pt,
  innertopmargin=8pt, innerbottommargin=8pt,
  innerleftmargin=10pt, innerrightmargin=10pt,
  roundcorner=4pt
}
\mdfdefinestyle{ok}{
  backgroundcolor={green!6},
  linecolor=vertOK, linewidth=1pt,
  innertopmargin=7pt, innerbottommargin=7pt,
  innerleftmargin=10pt, innerrightmargin=10pt,
  roundcorner=4pt
}
\mdfdefinestyle{warn}{
  backgroundcolor={orange!8},
  linecolor=orangeWARN, linewidth=1pt,
  innertopmargin=7pt, innerbottommargin=7pt,
  innerleftmargin=10pt, innerrightmargin=10pt,
  roundcorner=4pt
}
\mdfdefinestyle{danger}{
  backgroundcolor={red!6},
  linecolor=rougeERR, linewidth=1pt,
  innertopmargin=7pt, innerbottommargin=7pt,
  innerleftmargin=10pt, innerrightmargin=10pt,
  roundcorner=4pt
}

\lstset{
  basicstyle=\ttfamily\small,
  backgroundcolor=\color{grisCode},
  frame=single, rulecolor=\color{bleuCMR!30},
  breaklines=true, showstringspaces=false,
  tabsize=4, xleftmargin=4pt, xrightmargin=4pt
}

\pagestyle{fancy}
\fancyhf{}
\fancyhead[L]{\small\textcolor{bleuCMR}{Simulation de pannes — Hôpital camerounais}}
\fancyhead[R]{\small\textcolor{bleuCMR}{\thepage}}
\renewcommand{\headrulewidth}{0.4pt}

\newcommand{\ok}[1]{\textcolor{vertOK}{\textbf{#1}}}
\newcommand{\warn}[1]{\textcolor{orangeWARN}{\textbf{#1}}}
\newcommand{\err}[1]{\textcolor{rougeERR}{\textbf{#1}}}

% ═══════════════════════════════════════════════════════════════════════
\begin{document}
% ═══════════════════════════════════════════════════════════════════════

% ── Page de titre ──────────────────────────────────────────────────────
\begin{titlepage}
  \centering
  \vspace*{1.5cm}
  {\color{bleuCMR}\rule{\linewidth}{2pt}}
  \vspace{0.4cm}

  {\LARGE\bfseries Dispositif de prédiction de pannes d'électricité\\
  et de basculement automatique}

  \vspace{0.3cm}
  {\color{bleuCMR}\rule{\linewidth}{0.8pt}}
  \vspace{1cm}

  {\large\bfseries Rapport technique — Phase 1, Version 3}\\
  \vspace{0.3cm}
  {\large Simulation de capteurs ROS 2, collecte de données\\
  et analyse du dataset — Hôpital camerounais}

  \vspace{1.5cm}
  \begin{mdframed}[style=encadre]
    \centering
    \textbf{Environnement cible :} Hôpital $\sim$100 lits, Cameroun\\
    \textbf{Réseau simulé :} BTA 220 V / 50 Hz (Règlement MINEE art.\ 14)\\
    \textbf{Outil :} ROS 2 Jazzy (Python 3) — \texttt{sensor\_publisher v3.0}\\
    \textbf{Référence :} Moukengué Imano A., Université de Douala, 2015\\
    \textbf{Statut :} 15 458 lignes collectées — 3 sessions validées
  \end{mdframed}

  \vfill
  {\small Version 3.0 — \today}
\end{titlepage}

% ── Table des matières ─────────────────────────────────────────────────
\tableofcontents
\newpage

% ═══════════════════════════════════════════════════════════════════════
\section{Évolution du projet — de la v1.0 à la v3.0}
% ═══════════════════════════════════════════════════════════════════════

Ce rapport est la troisième version du document technique du projet.
Il intègre l'ensemble des corrections apportées au simulateur, les
résultats des premières sessions de collecte réelle, et les décisions
méthodologiques prises pour la construction du dataset d'entraînement.

\subsection{Historique des versions}

\begin{center}
\begin{tabular}{p{0.08\linewidth} p{0.22\linewidth} p{0.58\linewidth}}
\toprule
\textbf{Version} & \textbf{Date} & \textbf{Contenu principal} \\
\midrule
v1.0 & Phase initiale &
  Architecture ROS 2 définie, premiers fichiers Python créés,
  paramètres génériques non calés sur le contexte camerounais. \\
v2.0 & Après analyse manuscrit &
  Paramètres corrigés selon les Tableaux 1 et 2 du manuscrit
  (tension, fréquence), ajout THD et microcoupures,
  18 colonnes CSV définies. \\
\textbf{v3.0} & \textbf{Après validation datasets} &
  \textbf{5 bugs critiques corrigés, probabilités recalibrées,
  machine à états précurseur implémentée, 4 sessions de collecte
  analysées et validées.} \\
\bottomrule
\end{tabular}
\end{center}

% ═══════════════════════════════════════════════════════════════════════
\section{Corrections apportées dans la version 3.0}
\label{sec:corrections}
% ═══════════════════════════════════════════════════════════════════════

L'analyse des premiers datasets (sessions du 5 et 6 avril 2026) a
révélé cinq bugs critiques dans le code \texttt{sensor\_publisher.py}.
Ces bugs produisaient des données irréalistes rendant le dataset
inutilisable pour l'entraînement d'un modèle ML.

\subsection{BUG-1 — Probabilités de coupure trop élevées}

\begin{mdframed}[style=danger]
\textbf{Symptôme :} La session 1 (7h40) montrait 99.8\,\% de
coupures — une seule coupure de 157 minutes couvrait toute la
simulation. La session 2 (12h27) : 98.5\,\% de coupures.
\end{mdframed}

\textbf{Cause :} Les probabilités de déclenchement ($p = 2.5$ à
$5\,\%$ par cycle de 2\,s) étaient trop élevées.
À $p = 5\,\%$, l'espérance d'une coupure toutes les $1/(0{,}05)
\times 2 = 40$\,s rendait le réseau quasi-constamment en panne.

\textbf{Correction :} Les probabilités ont été recalibrées par
simulation statistique sur 20 runs de 1 heure afin d'obtenir
$\sim$92\,\% d'état normal en heure de pointe~:

\begin{center}
\begin{tabular}{lccc}
\toprule
\textbf{Plage horaire} & \textbf{Avant} & \textbf{Après} &
\textbf{\% normal attendu} \\
\midrule
Nuit (22h–06h)   & 1{,}0 \%  & 0{,}005 \%  & $\sim$99{,}8 \% \\
Journée (06h–18h) & 2{,}5 \%  & 0{,}015 \%  & $\sim$97{,}0 \% \\
Pointe (18h–22h) & 5{,}0 \%  & 0{,}030 \%  & $\sim$91{,}9 \% \\
\bottomrule
\end{tabular}
\end{center}

\subsection{BUG-2 — Durée maximale de coupure irréaliste}

\begin{mdframed}[style=danger]
\textbf{Symptôme :} Coupure unique de 157 minutes dans la session 1,
causée par un tirage aléatoire dans la plage 30–10 800 cycles.
\end{mdframed}

\textbf{Correction :} La durée maximale est réduite à
$450$ cycles ($= 15$\,min), cohérente avec les interventions
rapides documentées pour les réseaux urbains camerounais.

\begin{lstlisting}[language=Python]
# AVANT (6 heures max — irréaliste)
self.outage_duration = random.randint(30, 10_800)

# APRÈS (15 minutes max — réaliste)
self.outage_duration = random.randint(15, 450)
\end{lstlisting}

\subsection{BUG-3 — Puissance non nulle pendant une microcoupure}

\begin{mdframed}[style=danger]
\textbf{Symptôme :} Des lignes avec \texttt{micro\_outage=1},
\texttt{voltage=0V} mais \texttt{power=22 000 à 48 000 W}.
Physiquement impossible~: si la tension est nulle, la puissance
active l'est aussi ($P = U \times I$).
\end{mdframed}

\textbf{Correction :} La fonction \texttt{get\_power()} retourne
\texttt{0.0} dès que \texttt{micro\_outage\_active} est vrai,
exactement comme pour la coupure complète.

\subsection{BUG-4 — Fréquence de microcoupure trop élevée}

\textbf{Symptôme :} Avec $p = 4\,\%$ par cycle, une microcoupure
survenait toutes les 50 secondes en moyenne, rendant l'état
«~normal~» très rare.

\textbf{Correction :} Probabilité ramenée à $p = 0{,}2\,\%$,
produisant $\sim$3 à 5 microcoupures par heure — valeur
cohérente avec les démarrages de moteurs documentés au
§3.3.1 du manuscrit.

\subsection{BUG-5 — Signal précurseur absent}

\begin{mdframed}[style=warn]
\textbf{Contexte :} Le signal précurseur est la feature ML la
plus importante du projet. Il correspond à la chute progressive
de tension observée quelques secondes avant une coupure, décrite
au §3.3.1 du manuscrit comme «~creux de tension~».
\end{mdframed}

\textbf{Symptôme :} Dans v2.0, la chute de tension avant coupure
était implémentée via une condition mal placée qui ne se
déclenchait jamais.

\textbf{Correction :} Un état dédié \texttt{pre\_fault\_active}
a été ajouté à la machine à états. Pendant cet état (6 à 16\,s
avant la coupure), la tension chute progressivement de 0 à
$-50$\,V selon la progression du précurseur~:

\begin{lstlisting}[language=Python]
# Chute progressive proportionnelle à l'avancement
progress = 1.0 - (self.pre_fault_duration / self.pre_fault_total)
pre_drop = -random.uniform(10.0, 50.0) * progress
voltage  = VOLTAGE_NOMINAL + noise + pre_drop
\end{lstlisting}

\subsection{Résumé des corrections — tableau comparatif}

\begin{center}
\begin{tabularx}{\linewidth}{p{0.18\linewidth} p{0.25\linewidth}
                              p{0.25\linewidth} X}
\toprule
\textbf{Bug} & \textbf{Symptôme} & \textbf{Correction} &
\textbf{Résultat} \\
\midrule
BUG-1 proba &
  99\,\% coupures &
  $\div 100$ les probabilités &
  $\sim$92\,\% normal en pointe \\
BUG-2 durée &
  Coupure 157\,min &
  Max 15\,min (450 cycles) &
  Durées réalistes 1–15\,min \\
BUG-3 power &
  P $\neq$ 0 si V = 0 &
  \texttt{return 0.0} micro &
  Cohérence physique \\
BUG-4 micro &
  1 micro / 50\,s &
  $p : 4\,\% \to 0{,}2\,\%$ &
  3–5 micros par heure \\
BUG-5 précurseur &
  Jamais déclenché &
  État \texttt{pre\_fault} dédié &
  Chute visible 6–16\,s avant \\
\bottomrule
\end{tabularx}
\end{center}

% ═══════════════════════════════════════════════════════════════════════
\section{Machine à états du simulateur v3.0}
% ═══════════════════════════════════════════════════════════════════════

La version 3.0 implémente une machine à états à 4 niveaux de priorité,
exécutée à chaque cycle de 2 secondes~:

\begin{center}
\begin{tabular}{clp{0.55\linewidth}}
\toprule
\textbf{Priorité} & \textbf{État} & \textbf{Description} \\
\midrule
1 & \texttt{outage\_active}     &
  Tension, fréquence, puissance, THD = 0.
  Dure 30\,s à 15\,min. \\
2 & \texttt{micro\_outage\_active} &
  Idem mais 2 à 10\,s.
  Indépendant de la coupure complète. \\
3 & \texttt{pre\_fault\_active}  &
  Tension chute progressivement. Dure 6 à 16\,s.
  \textbf{C'est le signal que le ML doit détecter.} \\
4 & Normal                       &
  Mesures réalistes selon l'heure et le profil hôpital. \\
\bottomrule
\end{tabular}
\end{center}

\begin{mdframed}[style=encadre]
\textbf{Séquence type dans le CSV} (cycles consécutifs)~:\\[4pt]
\texttt{cycle n   : V=218V  outage=0  pre\_fault=0  (normal)}\\
\texttt{cycle n+1 : V=210V  outage=0  pre\_fault=0  (normal)}\\
\texttt{cycle n+2 : V=195V  outage=0  pre\_fault=1  (précurseur débute)}\\
\texttt{cycle n+3 : V=178V  outage=0  pre\_fault=1  (chute s'accentue)}\\
\texttt{cycle n+4 : V=155V  outage=0  pre\_fault=1  (chute forte)}\\
\texttt{cycle n+5 : V=0V    outage=1  pre\_fault=0  (COUPURE)}\\
\texttt{cycle n+6 : V=0V    outage=1  pre\_fault=0  (...)}\\
\texttt{cycle n+k : V=220V  outage=0  pre\_fault=0  (rétabli)}
\end{mdframed}

% ═══════════════════════════════════════════════════════════════════════
\section{Résultats des sessions de collecte}
\label{sec:sessions}
% ═══════════════════════════════════════════════════════════════════════

Quatre sessions de collecte ont été réalisées après la correction des
bugs. Trois d'entre elles ont été validées et retenues.

\subsection{Tableau de synthèse des sessions}

\begin{center}
\begin{tabularx}{\linewidth}{p{0.14\linewidth} p{0.11\linewidth}
                              p{0.08\linewidth} p{0.10\linewidth}
                              p{0.10\linewidth} p{0.10\linewidth} X}
\toprule
\textbf{Session} & \textbf{Plage} & \textbf{Lignes} &
\textbf{Normal} & \textbf{Coupures} & \textbf{Micros} &
\textbf{Statut} \\
\midrule
Nuit (v3.0)     & 00h–03h & 6 175  & 99{,}5\,\% & 0  & 10 &
  \ok{Valide} \\
Journée (v3.0)  & 13h–14h & 2 662  & 92{,}9\,\% & 1  &  6 &
  \ok{Valide} \\
Soirée (v3.0)   & 18h–21h & 6 621  & 95{,}9\,\% & 2  & 14 &
  \ok{Valide} \\
\midrule
Session 1 (v2.0) & 07h–10h & 4 731  & 0{,}2\,\%  & 1  &  0 &
  \err{Rejetée (BUG-2)} \\
Session 2 (v2.0) & 12h–12h & 995   & 0{,}9\,\%  & 1  &  0 &
  \err{Rejetée (BUG-1)} \\
\midrule
\multicolumn{2}{l}{\textbf{Total validé}} & \textbf{15 458} &
  \multicolumn{2}{l}{\textbf{3 coupures}} &
  \multicolumn{2}{l}{\textbf{30 microcoupures}} \\
\bottomrule
\end{tabularx}
\end{center}

\subsection{Évolution des mesures selon la plage horaire}

La cohérence des données entre sessions confirme que la simulation
reproduit fidèlement les patterns du réseau camerounais~:

\begin{center}
\begin{tabular}{lcccc}
\toprule
\textbf{Plage} & \textbf{V moy (V)} & \textbf{P moy (kW)} &
\textbf{THD moy (\%)} & \textbf{Coupures/h} \\
\midrule
Nuit (00h–03h)  & 221{,}2 & 105  & 12{,}2 & 0{,}00 \\
Journée (13h–14h) & 212{,}9 & 185  & 15{,}1 & 0{,}68 \\
Soirée (18h–21h) & 221{,}1 & 145  & 13{,}6 & 0{,}54 \\
\bottomrule
\end{tabular}
\end{center}

Ces évolutions sont physiquement cohérentes avec le manuscrit~:
\begin{itemize}
  \item La tension est plus basse en journée (réseau surchargé,
    §3.3.1), conforme au Tableau~1 (moyenne mesurée 213{,}4\,V).
  \item La puissance atteint son maximum en journée (blocs
    opératoires actifs, 185\,kW = 92\,\% du nominal).
  \item Les coupures surviennent uniquement aux heures de pointe
    (19h et 21h), fidèle au §3.4.
\end{itemize}

\subsection{Validation de la conformité légale}

\begin{center}
\begin{tabular}{lccc}
\toprule
\textbf{Plage} &
\textbf{V hors ±10\,\%} &
\textbf{f hors ±5\,\%} &
\textbf{Cohérence physique} \\
\midrule
Nuit     & 3{,}9\,\%  & 23{,}6\,\% & \ok{OK} \\
Journée  & 13{,}9\,\% & 44{,}3\,\% & \ok{OK} \\
Soirée   & 4{,}0\,\%  & 23{,}3\,\% & \ok{OK} \\
\bottomrule
\end{tabular}
\end{center}

\begin{mdframed}[style=warn]
\textbf{Note sur la fréquence hors tolérance~:} le taux de 44\,\%
en journée est élevé mais \textbf{non pathologique}. Le Tableau~2
du manuscrit documente des fréquences jusqu'à 66–70\,Hz mesurées
à Douala en journée. Ces valeurs extrêmes rares font monter le
pourcentage statistique. Ce signal est une \textbf{feature
discriminante utile} pour le modèle ML, pas une erreur.
\end{mdframed}

% ═══════════════════════════════════════════════════════════════════════
\section{Analyse du déséquilibre de classes}
\label{sec:desequilibre}
% ═══════════════════════════════════════════════════════════════════════

\subsection{Le problème fondamental}

L'analyse des datasets révèle un déséquilibre structurel inhérent
au projet~: les coupures sont des événements \textbf{rares par nature}.

\begin{center}
\begin{tabular}{lcc}
\toprule
\textbf{Classe} & \textbf{Lignes sur 15 458} & \textbf{Proportion} \\
\midrule
Normal (\texttt{outage = 0}) & $\sim$15 200 & $\sim$98{,}4\,\% \\
Coupure (\texttt{outage = 1}) & $\sim$242   & $\sim$1{,}6\,\% \\
\bottomrule
\end{tabular}
\end{center}

Un modèle ML entraîné sur un tel dataset apprendra à prédire
\textbf{toujours «~normal~»} et obtiendra une précision
apparente de 98\,\% — mais sera \textbf{totalement inutile} pour
détecter les vraies pannes.

\subsection{Pourquoi la simulation pure ne suffit pas}

À raison de 0{,}54 coupure par heure en heure de pointe, il faudrait
environ \textbf{620 heures de simulation} pour atteindre 500 coupures.
Ce n'est pas faisable dans le calendrier du projet.

\subsection{La solution : SMOTE}

\textbf{SMOTE} (\textit{Synthetic Minority Over-sampling Technique})
est la technique standard pour corriger le déséquilibre de classes.
Elle génère des exemples synthétiques de la classe minoritaire par
interpolation entre des exemples réels existants — sans inventer
de données incohérentes.

\begin{mdframed}[style=ok]
\textbf{Plan d'action adopté~:}
\begin{enumerate}[leftmargin=*, nosep]
  \item Continuer la collecte jusqu'à \textbf{30 à 50 coupures réelles}
    (priorité aux sessions 18h–22h).
  \item Fusionner tous les CSV validés en un seul dataset.
  \item Appliquer SMOTE sur la classe \texttt{outage = 1}
    pour atteindre $\sim$10\,\% de coupures.
  \item Entraîner le modèle sur le dataset rééquilibré.
\end{enumerate}
\end{mdframed}

% ═══════════════════════════════════════════════════════════════════════
\section{État d'avancement et prochaines étapes}
% ═══════════════════════════════════════════════════════════════════════

\subsection{Progression vers les objectifs}

\begin{center}
\begin{tabular}{lccc}
\toprule
\textbf{Objectif} & \textbf{Cible} & \textbf{Actuel} &
\textbf{Progression} \\
\midrule
Lignes totales       & 50 000  & 15 458 & 31\,\%  \\
Coupures complètes   & $\geq 30$ & 3    & 10\,\%  \\
Heures simulées (6h–18h) & couverture & partielle & $\sim$10\,\% \\
Heures simulées (18h–22h) & couverture & 3h40 & $\sim$92\,\% \\
Heures simulées (22h–06h) & couverture & 3h26 & $\sim$43\,\% \\
\bottomrule
\end{tabular}
\end{center}

\subsection{Sessions prioritaires restantes}

\begin{center}
\begin{tabular}{clp{0.45\linewidth}}
\toprule
\textbf{Priorité} & \textbf{Plage horaire} & \textbf{Justification} \\
\midrule
1 & 08h–12h (journée) &
  Heures de montée en charge hôpital, blocs opératoires.
  Sous-représentées dans le dataset actuel. \\
2 & 18h–22h (pointe) &
  Maximum de coupures attendu. Déjà 3h40 mais insuffisant. \\
3 & 05h–07h (transition) &
  Passage nuit/journée, variations importantes de charge. \\
\bottomrule
\end{tabular}
\end{center}

% ═══════════════════════════════════════════════════════════════════════
\section*{Conclusion}
\addcontentsline{toc}{section}{Conclusion}
% ═══════════════════════════════════════════════════════════════════════

La version 3.0 du simulateur produit des données \textbf{correctes,
cohérentes et fidèles} aux mesures réelles du réseau camerounais
documentées par Moukengué Imano. Les cinq bugs identifiés ont été
entièrement corrigés et les trois sessions de collecte validées
confirment que le pipeline ROS~2 fonctionne de bout en bout.

Le principal défi restant est le \textbf{déséquilibre structurel
des classes} : les coupures sont des événements rares ($\sim 1{,}6\,\%$
du temps), ce qui est réaliste mais incompatible avec un entraînement
ML direct. La stratégie SMOTE permettra de résoudre ce problème sans
dégrader la qualité des données.

Le pipeline complet — collecte $\to$ fusion $\to$ SMOTE $\to$ modèle
$\to$ basculement — sera documenté dans le rapport de Phase~2.

% ── Références ─────────────────────────────────────────────────────────
\begin{thebibliography}{9}

\bibitem{moukengue2015}
  Moukengué Imano A.,
  \textit{Problèmes d'électrification urbaine au Cameroun :
  diagnostic et proposition de solutions curatives},
  Équipe de Recherche en Systèmes d'Énergie Électrique,
  Laboratoire LEEAT, Université de Douala, Novembre 2015.

\bibitem{minee2009}
  Ministère de l'Énergie et de l'Eau du Cameroun,
  \textit{Arrêté n°~00000013/MINEE du 26~janvier~2009 portant
  approbation du Règlement du Service de distribution publique
  d'électricité}, Yaoundé, 2009.

\bibitem{loi2011}
  République du Cameroun,
  \textit{Loi N°~2011/022 du 14~décembre~2011 régissant le secteur
  d'électricité au Cameroun}, Assemblée Nationale, Yaoundé, 2011.

\bibitem{chawla2002}
  Chawla N.V., Bowyer K.W., Hall L.O., Kegelmeyer W.P.,
  \textit{SMOTE: Synthetic Minority Over-sampling Technique},
  Journal of Artificial Intelligence Research, vol.\ 16,
  pp.\ 321--357, 2002.

\end{thebibliography}

\end{document}
